\documentclass{article}
\begin{document}

\title{Asset Viewer: LaTeX Sample}
\author{Barodoc Team}
\maketitle

\section{Introduction}

This sample is written for the \textbf{LaTeX.js} viewer in the browser. It uses only the default \texttt{article} class and no extra packages (\verb|\usepackage|). Packages like \texttt{babel}, \texttt{geometry}, or \texttt{amsmath} are loaded by LaTeX.js with \texttt{require()}, which is not available in the browser, so they are omitted here.

\section{Structure}

\subsection{Sections and text}

You can use \textbf{bold}, \emph{emphasis}, and \texttt{code} style. Sections and subsections are rendered as headings.

\subsection{List}

\begin{itemize}
  \item First item.
  \item Second item.
\end{itemize}

\begin{enumerate}
  \item First numbered item.
  \item Second numbered item.
\end{enumerate}

\subsection{Table}

Browser LaTeX.js does not support the \texttt{table} environment; use \texttt{tabular} only. Example:

\begin{center}
  \begin{tabular}{|l|l|}
    \hline
    Format & Viewer \\
    \hline
    .tex & LaTeX.js \\
    .csv & Table \\
    .rst & RST parser \\
    \hline
  \end{tabular}
\end{center}

\section{Conclusion}

If you see this as formatted text with title, sections, list, and table, the LaTeX viewer is working. For documents that need \texttt{amsmath}, \texttt{babel}, or \texttt{geometry}, use a build-time conversion (e.g. \texttt{pdflatex} or LaTeX.js CLI) and serve the generated HTML or PDF instead.

\end{document}
